\section{Two seeds with different orders}\label{seed:section}

We first construct two sums of chains with the same spectrum modulo~$2$,
but with ordinary orders $2$ and $3$.  The arithmetic ingredient is a
finite relation between positive rational numbers.

\begin{lemma}[Finite product relation]\label{seed:product}
Set $d=3^{11}$ and
\[
 \mathcal B=\{b\in\Z:0<b<d/2,\ 3\nmid b\},\qquad
 A_b=\frac{d-b}{b},\qquad
 R_b=\frac{(2d+b)(2d-b)}{(3d-b)(d+b)}.
\]
There are integers $e_b$, $b\in\mathcal B$, such that
\begin{equation}\label{seed:products}
 E:=\sum_{b\in\mathcal B}e_b>0,
 \qquad \prod_{b\in\mathcal B}A_b^{e_b}
 =\prod_{b\in\mathcal B}R_b^{e_b}=1.
\end{equation}
\end{lemma}

\begin{proof}
Form the integer matrix with column $b$ consisting of
$(v_p(A_b))_p$ and $(v_p(R_b))_p$, with separate rows for the two
families of prime valuations.  It is enough to find an integer kernel
vector whose coordinate sum is nonzero.

Here is the finite certificate for this assertion.  Retain only the
columns for which each of
\[
 d-b,\ b,\ 2d+b,\ 2d-b,\ 3d-b,\ d+b
\]
has no prime factor greater than $4450$.  Repeatedly delete any column
that is the only column supported on some nonzero row; discard zero
rows.  The resulting matrix $M$ and its augmentation $\widetilde M$
by a final row of ones satisfy
\begin{center}
\begin{tabular}{lrr}
 & columns & nonzero valuation rows\\\hline
After the prime cutoff & $1118$ & $1100$\\
After all deletions & $896$ & $867$
\end{tabular}
\end{center}
and
\begin{equation}\label{seed:rank-certificate}
 \operatorname{rank}_{\F_{1009}}\widetilde M=868.
\end{equation}
Appendix~\ref{app:certificate} gives a complete program, using only
integer arithmetic and elimination modulo the prime $1009$, that
verifies these counts and~\eqref{seed:rank-certificate}.  Its sieve
stores a prime divisor of every integer greater than one.  Its six
signed factor entries are precisely the two valuation columns above.
In the elimination, a normalized pivot is stored for each leading row,
so the number of stored pivots is the rank.  Thus the certificate is a
finite check of the displayed integer matrix, without numerical
approximation.

Since $M$ has $867$ rows,~\eqref{seed:rank-certificate} gives an
$868$-column square submatrix $H$ of $\widetilde M$ whose integer
determinant is nonzero.  The vector
\[
 z=\operatorname{adj}(H)(0,\ldots,0,1)^{\mathsf t}
\]
is integral, satisfies the valuation equations on these columns, and
has coordinate sum $\det H$.  Extend $z$ by zeros to all columns
$b\in\mathcal B$, and change its overall sign if necessary.  This
gives $e_b$ with $E=|\det H|>0$.  Every omitted valuation row is zero on
the retained columns, so all prime valuations of both products
in~\eqref{seed:products} vanish.  A positive rational number with this
property is $1$, proving the Lemma.
\end{proof}

The adjugate prescription also makes the construction entirely finite:
one may select the lexicographically first nonsingular minor, with
columns ordered by $b$ and valuation rows ordered by family and prime.
No bound on the size of its coefficients is needed.

\begin{lemma}[Completing one weight to a chain]\label{seed:chain}
Let $D$ be a positive odd integer and let $c\in\Z$ satisfy
$0<c<D/2$.  There is a chain
polynomial
\begin{equation}\label{seed:chain-polynomial}
 F_c=x_1^{a_1}x_2+\cdots+x_{\ell-1}^{a_{\ell-1}}x_\ell
       +x_\ell^{a_\ell}
\end{equation}
of weighted degree $1$, whose first weight is $c/D$, all of whose
weights are positive and less than $1/2$, and all of whose ordinary
monomial degrees are at least $3$.  It has an isolated critical point
at the origin.  Removing its first variable and first monomial gives
a chain with exactly the remaining weights, when $\ell>1$.
\end{lemma}

\begin{proof}
Set $c_1=c$ and successively perform the divisions
\[
 a_i=\lfloor D/c_i\rfloor,\qquad c_{i+1}=D-a_ic_i,
\]
stopping at the first zero remainder.  Each nonzero remainder is a
positive integer strictly smaller than its predecessor, so this
procedure terminates.
Give $x_i$ weight $c_i/D$.  Then $a_ic_i+c_{i+1}=D$ for each
nonterminal term, and $a_\ell c_\ell=D$ for the terminal term.
All $a_i\ge2$.  Consequently each nonterminal monomial has ordinary
degree $a_i+1\ge3$.  The terminal exponent divides the odd integer
$D$ and exceeds $2$, so it too is at least $3$.

To prove isolation, suppose a nonzero point is critical and let $j$
be the first index with $x_j\ne0$.  If $j=\ell$, the derivative in
$x_j$ is $a_\ell x_\ell^{a_\ell-1}\ne0$.  Otherwise that derivative
forces $x_{j+1}=0$, because the preceding coordinate is zero.  The
derivative in $x_{j+1}$ is then $x_j^{a_j}\ne0$: its other term
vanishes since $a_{j+1}\ge2$.  Both cases contradict criticality.
The same reasoning applies to the tail, which has the stated weights.
\end{proof}

\begin{proposition}[The seeds]\label{seed:seeds}
There are real-coefficient weighted homogeneous polynomials
$p_1,p_2$, each in an even number $s_0$ of variables, such that:
\begin{enumerate}
\item both are sums, in disjoint variables, of chains with exponents
      at least $2$, and both have isolated critical points;
\item $\ord_0p_1=2$ and $\ord_0p_2=3$;
\item their spectral multisets are equal modulo $2$;
\item their real Seifert forms are isomorphic, and, separately, their
      integral monodromy modules are isomorphic;
\item both monodromy operators have order dividing $D=3^{12}$.
\end{enumerate}
\end{proposition}

\begin{proof}
Fix the integers in Lemma~\ref{seed:product}, and put $D=3d$.
The four weights below are chosen for the triple-angle identity in the
spectral comparison; their signed counts will give side $2$ an excess
of $4E$ variables, to be balanced by squares on side $1$.
Start with two lists of weights prescribed by the following signed
counts; a positive entry means that many variables on side $2$, and a
negative entry means its absolute value on side $1$:
\begin{equation}\label{seed:weight-table}
\begin{array}{c|rrrr}
 w& b/d&b/(3d)&(d-b)/(3d)&(d+b)/(3d)\\\hline
 \text{count on side $2$ minus count on side $1$}
   &2e_b&-2e_b&2e_b&2e_b.
\end{array}
\end{equation}
All these weights are positive and less than $1/2$.  Their numerators
with denominator $D$ are $3b,b,d-b,d+b$, respectively.

Realize each occurrence of a starting weight $c/D$ by its own copy of
the chain in Lemma~\ref{seed:chain}.  On the opposite side insert an
independent copy of the tail of that chain, doing nothing there if the
tail is empty.  The added tail weights occur equally on both sides.
All variables used for different copies are disjoint.  Thus this
operation realizes the signed weight counts exactly by sums of
isolated chains, and all its monomials have degree at least $3$.

The variable excess on side $2$ is $4E$.  Add $4E$ independent squares
to side $1$, giving both sides the same number of variables.  Every
starting multiplicity in~\eqref{seed:weight-table} is even, as is the
number of copies of each tail, and $4E$ is even.  Both variable counts
are therefore even already.  Finally add two independent pure cubes
to each side, and call the resulting polynomials $p_1,p_2$.  They have
the same even number $s_0$ of variables.  Sums in disjoint variables
have critical loci equal to products of the individual critical loci,
so both critical points are isolated.  All coefficients can be taken
to be $1$.  Squares occur only in $p_1$, and the common cubes occur in
both polynomials, proving the asserted ordinary orders.

We next compare the spectra.  In the convention of
Theorem~\ref{bg:spectrum}, write
\begin{equation}\label{seed:spectrum-product}
 P_p(t)=\sum_{\alpha\in\operatorname{Sp}^{+}(p)}t^\alpha
        =\prod_{w\text{ a weight of }p}\frac{t^w-t}{1-t^w}.
\end{equation}
All exponents lie in $(2D)^{-1}\Z$.  Reduction modulo $2$ puts their
multiplicities in the group algebra of the finite cyclic group
$(2D)^{-1}\Z/2\Z$.  Its characters are
$\alpha\mapsto e^{\pi i k\alpha}$, for integer $k$.  By Fourier
inversion, it suffices to prove that $P_{p_1}$ and $P_{p_2}$ take the
same value at every such character.  When an individual quotient
in~\eqref{seed:spectrum-product} is indeterminate, evaluate along
$t^u=\exp((\pi i k+\varepsilon)u)$ and let $\varepsilon\to0$.

\smallskip\noindent
\emph{Case 1: $k=2m$ is even.}
For a weight $w$, its individual factor has value
\begin{equation}\label{seed:even-factor}
 \left.\frac{t^w-t}{1-t^w}\right|_k
 =\begin{cases}
 -1,&mw\notin\Z,\\
 (1-w)/w,&mw\in\Z.
 \end{cases}
\end{equation}
The second value follows by taking the ratio of the two first
derivatives in $\varepsilon$.  All these values are nonzero, so common
weights can be canceled.

The first column of~\eqref{seed:weight-table} has exact reduced
denominator $d$, and the other three have exact reduced denominator
$3d$, since $3\nmid b$.  Within either denominator group, therefore,
all factors are exceptional in~\eqref{seed:even-factor}, or none are.
The exceptional signed product for the first group is
$\prod_b A_b^{2e_b}=1$.  For the other group it is
\[
 \prod_b\left(
 \frac{b}{3d-b}\frac{2d+b}{d-b}\frac{2d-b}{d+b}
 \right)^{2e_b}
 =\prod_b(R_b/A_b)^{2e_b}=1.
\]
The nonexceptional signs disappear because all signed counts are
even.  A square has exceptional value $1$ and nonexceptional value
$-1$; its total contribution is $1$, since there are $4E$ squares.
This proves equality in Case~1, including $k=0$.

\smallskip\noindent
\emph{Case 2: $k$ is odd and $3\nmid k$.}
Each factor is
\[
 \left.\frac{t^w-t}{1-t^w}\right|_k
     =i\cot(\pi k w/2).
\]
All chain weights have a reduced denominator that is a nontrivial
power of $3$.  None of their cotangents has a zero or pole in this
case: either would imply that this denominator divides $k$.
Squares also have finite nonzero factors.  Hence common weights may
again be canceled.

Put $C(u)=\cot(\pi ku/2)$.  The exact signed identity we need is
\begin{equation}\label{seed:cotangent}
 C(b/d)=
 \frac{C(b/(3d))}{C((d-b)/(3d))C((d+b)/(3d))}.
\end{equation}
Indeed, with $x=\pi kb/(6d)$ the usual triple-angle product gives
\[
 \cot x\,\cot(x+k\pi/3)\,\cot(x+2k\pi/3)=-\cot(3x).
\]
The shifts by $k\pi/3$ give the same three factors modulo $\pi$ as
the shifts by $\pi/3$, because $3\nmid k$.  Since $k$ is odd,
\[
 \cot(x+k\pi/3)=\frac{1}{C((d-b)/(3d))},\qquad
 \cot(x+2k\pi/3)=-\frac{1}{C((d+b)/(3d))}.
\]
Substitution proves~\eqref{seed:cotangent}.  Its $2e_b$-th power is
exactly the assertion that the signed cotangent product in row $b$
of~\eqref{seed:weight-table} equals $1$.
Each square contributes $C(1/2)=\cot(\pi k/4)\in\{1,-1\}$, whose
$4E$-th power is $1$.  Finally the powers of $i$ are the same because
the two complete polynomials have the same number $s_0$ of variables.
This proves equality in Case~2.

\smallskip\noindent
\emph{Case 3: $k$ is odd and $3\mid k$.}
A pure cube has spectrum $t^{1/3}+t^{2/3}$, whose value is $-1+1=0$.
Before the final common cubes were inserted, both sides were already
isolated sums of chains and squares, so their spectra were finite
polynomials.  Multiplication by a cube spectrum therefore makes both
character values zero.  This argument uses those finite polynomials,
and does not cancel a zero against a pole of an individual weight
quotient.  It completes the Fourier comparison.

Weighted homogeneous monodromy has finite order and hence is
semisimple.  Theorem~\ref{bg:real} now converts the equality of
spectra modulo $2$ in the common dimension into an isomorphism of
real Seifert forms.  The same equality gives equal eigenvalue
multiplicities.  All summands are chains, including the pure powers,
so Theorem~\ref{bg:orlik} gives an isomorphism of integral monodromy
modules.  These are separate isomorphisms; no compatibility between
them is asserted here.

Finally, every nonsquare weight has denominator dividing $D$.
Each square has spectrum $t^{1/2}$, so the $4E$ squares together
contribute $t^{2E}$.  All spectral exponents of either seed therefore
belong to $D^{-1}\Z$.  Their monodromy eigenvalues are $D$-th roots of
unity.  Semisimplicity then gives $M_{p_j}^D=1$ for $j=1,2$, as claimed.
\end{proof}

\begin{corollary}[Removing eigenvalue one]\label{seed:prepared}
In two new variables on each side, put
\begin{equation}\label{seed:q-definition}
 q_j=p_j+x^5+y^7,\qquad j=1,2,
 \qquad s=s_0+2.
\end{equation}
These are real-coefficient isolated sums of chains in the same even
number $s$ of variables, with orders $2$ and $3$.  Their real Seifert
forms are isomorphic, and their integral monodromy modules are
isomorphic.  Both monodromies have odd order dividing
\[
 m=3^{12}\cdot5\cdot7,
\]
and neither has eigenvalue $1$.
\end{corollary}

\begin{proof}
The new terms preserve isolation and ordinary orders.  Formula
\eqref{seed:spectrum-product} multiplies each seed spectrum by the
same two pure-power spectra, preserving equality modulo $2$.
Theorems~\ref{bg:real} and~\ref{bg:orlik} consequently apply again.
Each new eigenvalue has the form
\[
 \alpha\beta\gamma,\qquad
 \alpha^D=1,\quad
 \beta^5=\gamma^7=1,\quad \beta\ne1,\quad\gamma\ne1.
\]
Because $D$, $5$, and $7$ are pairwise coprime, its order divides
$35D$ and is divisible by both $5$ and $7$.  In particular it is not
$1$.  Finite-order weighted homogeneous monodromy is semisimple, so
the same bound holds for the order of the operators themselves.
\end{proof}
